%% Springer LNCS / ECIR / CIKM template
%% Compile with: pdflatex main.tex && bibtex main && pdflatex main.tex && pdflatex main.tex

\documentclass[runningheads]{llncs}

% ── Packages ──────────────────────────────────────────────────────────────────
\usepackage[T1]{fontenc}
\usepackage[utf8]{inputenc}
\usepackage{graphicx}
\usepackage{hyperref}
\usepackage{color}
\usepackage{booktabs}
\usepackage{multirow}
\usepackage{amsmath}
\usepackage{amssymb}
\usepackage{enumitem}
\usepackage{xcolor}
\usepackage{microtype}
\renewcommand\UrlFont{\color{blue}\rmfamily}

% ── Metadata ──────────────────────────────────────────────────────────────────
\begin{document}

\title{Can You Tell the Difference? \\
  LLM-as-Discriminator vs.\ Human Judgment \\
  for Real and Synthetic Tabular Data}

\titlerunning{LLM-as-Discriminator for Tabular Data}

\author{Anonymous Author(s)}
\authorrunning{Anon.}
\institute{Vrije Universiteit Amsterdam, Amsterdam, The Netherlands}

\maketitle

% ── Abstract ──────────────────────────────────────────────────────────────────
\begin{abstract}
The growing adoption of synthetic tabular data generation has made
reliable evaluation of data fidelity increasingly critical.
Existing metrics---such as Total Variation Distance or Kolmogorov--Smirnov
statistics---operate at the aggregate level and do not capture the
\emph{perceived} realism of individual records or tables.
In this paper, we propose a \emph{LLM-as-Discriminator} framework in
which large language models (LLMs) are tasked with distinguishing real
from synthetic tabular data, and we compare their discriminative ability
against human judgments collected through a controlled labeling interface.
We investigate two information conditions: \textbf{C1}~(raw table only)
and \textbf{C2}~(raw table with distributional metadata including
statistics, correlations, and normality tests).
Our study examines whether LLMs can serve as scalable, cost-effective
proxies for human perception of data fidelity, and whether metadata
availability improves discrimination for both evaluator types.
\keywords{Synthetic tabular data \and Fidelity evaluation \and
LLM-as-evaluator \and Human evaluation \and Data generation}
\end{abstract}

% ── 1. Introduction ───────────────────────────────────────────────────────────
\section{Introduction}
\label{sec:intro}

Tabular data is the predominant format in institutional, clinical, and
financial settings, yet it is frequently subject to privacy constraints
that limit sharing and reuse~\cite{jordon2022synthetic}.
Synthetic data generation has emerged as a practical solution:
models such as CTGAN~\cite{xu2019modeling}, TVAE, GaussianCopula,
and more recently language-model-based generators such as
GReaT~\cite{borisov2023language} can produce tables that statistically
resemble their real counterparts.

Despite this progress, evaluating the \emph{fidelity} of synthetic
tabular data remains an open challenge.
Current evaluation protocols rely on aggregate metrics---distributional
divergences, pairwise correlations, or machine-learning utility
scores~\cite{platzer2021holdout}---which obscure local anomalies and
semantic implausibilities at the record level.
The question of whether a \emph{human expert} or an \emph{LLM} can
reliably detect that a table has been synthetically generated has
received virtually no attention for tabular data, despite being
extensively studied in image~\cite{rossler2019faceforensics} and
text~\cite{ippolito2020automatic} domains.

We address this gap by framing the discrimination task explicitly:
given a table, decide whether it is \textbf{REAL} (drawn from actual
observations) or \textbf{SYNTHETIC} (machine-generated).
We compare two classes of evaluators---LLMs and humans---across two
information conditions, yielding a $2 \times 2$ factorial design.

\paragraph{Contributions.}
\begin{enumerate}[leftmargin=*,label=(\roman*)]
  \item We introduce a \emph{LLM-as-Discriminator} protocol for tabular
        data, with structured prompts eliciting verdicts, confidence
        scores, and natural-language reasoning.
  \item We design and release a human labeling interface (Streamlit)
        that collects discrimination judgments with participant-level
        confidence and open-ended reasoning.
  \item We conduct a controlled experiment across two information
        conditions (raw table vs.\ table + metadata) and two evaluator
        types (LLM vs.\ human), reporting accuracy, calibration, and
        the qualitative cues each evaluator type relies upon.
  \item We provide an open-source framework\footnote{Code available at
        \url{https://github.com/anonymous/llm-discriminator}.}
        for replication and extension.
\end{enumerate}

% ── 2. Related Work ───────────────────────────────────────────────────────────
\section{Related Work}
\label{sec:related}

\subsection{Synthetic Tabular Data Generation}

The landscape of synthetic tabular data generation has evolved rapidly.
The \emph{Synthetic Data Vault} (SDV)~\cite{patki2016synthetic} provided
one of the first unified frameworks for multi-table synthesis using
Gaussian copulas and parametric models.
\citet{xu2019modeling} introduced CTGAN, which adapts the conditional
generative adversarial network paradigm to handle the mixed-type,
sparse, and imbalanced nature of real tabular data through
mode-specific normalisation of continuous columns and conditional
vector training.
Concurrent work on TVAE~\cite{xu2019modeling} applied variational
autoencoders to the same setting.

More recently, large language models have been explored as
tabular generators.
\citet{borisov2023language} demonstrated that GPT-2 fine-tuned on
serialised rows (\emph{GReaT}: \textit{Generation of Realistic
Tabular data}) can outperform GAN-based methods on structured fidelity
metrics while capturing complex inter-column dependencies.
REaLTabFormer~\cite{solatorio2023realtabformer} extended this approach
to relational tabular structures using encoder--decoder transformers.
ForestDiffusion~\cite{jolicoeur2023generating} applied score-based
diffusion models to tabular settings, leveraging tree-ensemble density
estimators to model heterogeneous column types.

\subsection{Evaluation of Synthetic Tabular Data}

Fidelity evaluation of synthetic tabular data can be broadly
classified into three paradigms~\cite{jordon2022synthetic}:

\begin{description}[leftmargin=0pt]
  \item[Statistical fidelity.] Metrics such as Total Variation Distance
    (TVD), Kolmogorov--Smirnov (KS) tests per column, and Wasserstein
    distance compare marginal and joint distributions between real and
    synthetic datasets~\cite{platzer2021holdout}.
    The SDMetrics library~\cite{datacebo2023sdmetrics} provides a
    comprehensive suite of such measures.

  \item[ML utility.] Train-on-synthetic, test-on-real (TSTR) and its
    inverse (TRTS) protocols assess whether a downstream ML model
    trained on synthetic data matches the performance of one trained on
    real data~\cite{esteban2017real}.

  \item[Privacy.] Membership inference, attribute inference, and
    nearest-neighbour distance metrics quantify the risk of re-identification
    from synthetic records~\cite{jordon2022synthetic}.
\end{description}

Critically, none of these paradigms assess the \emph{perceptual
realism} of individual tables---the question of whether a reasonably
informed observer can detect synthesis.
Our work introduces this fourth evaluation axis.

\subsection{LLM-as-Evaluator}

The use of LLMs as judges or evaluators has been extensively studied
in natural language processing.
\citet{zheng2023judging} introduced \emph{LLM-as-a-Judge} (MT-Bench),
demonstrating that GPT-4 judgments on open-ended generation tasks
correlate strongly with human preference ratings.
\citet{chiang2023large} evaluated whether LLMs can replace human
annotators in NLP benchmarks, finding high agreement for well-defined
tasks but systematic biases (e.g., position bias, verbosity preference).
\citet{wang2023large} proposed \emph{PandaLM}, a reproducible and
automated evaluation framework built around instruction-tuned LLMs.

For tabular data specifically, LLMs have been applied to table
question-answering~\cite{herzig2020tapas} and semantic parsing of
structured data~\cite{ye2023large}, but not to the problem of
discriminating real from synthetic tables.
Our work bridges this gap by adapting the LLM-as-evaluator paradigm
to a binary discrimination task over tabular inputs.

\subsection{Human Perception of Synthetic Data}

Human studies on synthetic data detection are common in vision
(deepfake detection~\cite{rossler2019faceforensics}) and
text~\cite{ippolito2020automatic,clark2021all}, but largely absent for
tabular data.
\citet{ippolito2020automatic} showed that human raters achieve only
modest accuracy at distinguishing human-written from GPT-2-generated
text, and that crowd-sourced accuracy is highly variable across domains.
\citet{clark2021all} conducted a large-scale human evaluation of
neural text generation quality, finding that humans systematically
conflate fluency with factuality.
These findings motivate us to study whether analogous confounds
arise in the tabular domain---for instance, whether humans mistake
statistically clean but semantically implausible rows as real.

% ── 3. Research Questions ─────────────────────────────────────────────────────
\section{Research Questions}
\label{sec:rqs}

Our study is organised around five research questions, reflecting a
$2 \times 2$ factorial design (\textbf{Evaluator type}: LLM vs.\ Human)
$\times$ (\textbf{Condition}: C1 vs.\ C2):

\begin{description}[leftmargin=0pt,itemsep=6pt]

  \item[\textbf{RQ1} --- Baseline Discriminability.]
    \emph{Can LLMs and humans distinguish real from synthetic tabular
    data at above-chance accuracy when presented with the raw table
    alone (Condition C1)?}

    This establishes whether the discrimination task is tractable at
    all for each evaluator type, and sets the baseline against which
    the effect of additional information is measured.

  \item[\textbf{RQ2} --- Effect of Metadata (C1 vs.\ C2).]
    \emph{Does providing distributional metadata (column statistics,
    pairwise correlations, and normality test results) significantly
    improve discrimination accuracy for LLMs and/or humans, relative
    to the raw-table-only condition?}

    We hypothesise that metadata systematically surfaces statistical
    artifacts of synthesis (e.g., irregular skewness, implausible
    inter-column correlations) that are difficult to detect from a
    visual scan of a table.

  \item[\textbf{RQ3} --- LLM vs.\ Human Comparison.]
    \emph{Do LLMs and humans differ in their overall discrimination
    accuracy, confidence calibration, and error patterns across
    conditions C1 and C2?}

    This RQ examines whether LLMs can serve as scalable proxies for
    human judgment in data fidelity assessment.
    We pay particular attention to asymmetric errors: whether each
    evaluator type disproportionately mislabels real data as synthetic
    (over-suspicious) or synthetic data as real (over-credulous).

  \item[\textbf{RQ4} --- Cue Utilisation.]
    \emph{Do LLMs and humans rely on qualitatively different cues when
    making discrimination judgments?}

    LLM responses include structured reasoning fields
    (\texttt{red\_flags}, \texttt{supporting\_evidence}) that
    allow cue analysis.
    Human participants provide open-ended justifications.
    We compare cue types across evaluator types and correctness
    outcomes to characterise \emph{what} makes synthetic tables
    detectable.

  \item[\textbf{RQ5} --- Correlation with Fidelity Metrics.]
    \emph{Does LLM discrimination accuracy correlate with
    established fidelity metrics (TVD, KS statistic, ML utility)
    across different synthetic datasets?}

    If harder-to-detect tables also score higher on standard metrics,
    LLM discrimination provides a new perceptual validity signal
    that can complement aggregate metrics.

\end{description}

% ── 4. Methodology ────────────────────────────────────────────────────────────
\section{Methodology}
\label{sec:method}

\subsection{Experimental Design}

We adopt a $2 \times 2$ factorial design crossing evaluator type
(LLM, Human) with information condition (C1, C2), yielding four
experimental cells.
Within each cell, trials are balanced across real and synthetic labels
and randomised in presentation order.

\paragraph{Condition C1 (Table only).}
The evaluator is presented with the first $N$ rows of a tabular
dataset (default $N = 20$), column names, and total dataset shape.
No statistical summary is provided.

\paragraph{Condition C2 (Table + Metadata).}
The evaluator receives the same table as in C1, augmented with a
structured metadata block containing: per-column descriptive statistics
(mean, standard deviation, min/max, quartiles, skewness, kurtosis),
a Shapiro--Wilk normality test result, top-$k$ value frequencies for
categorical columns, Shannon entropy, and pairwise Pearson correlations
among numeric columns.

\subsection{Datasets}

We evaluate on a suite of benchmark tabular datasets spanning diverse
domains and column types, paired with synthetic counterparts generated
by at least three synthesis methods:

\begin{table}[ht]
\centering
\caption{Datasets used in the study. \emph{Mixed} denotes datasets
containing both numeric and categorical columns.}
\label{tab:datasets}
\begin{tabular}{llrrl}
\toprule
\textbf{Dataset} & \textbf{Domain} & \textbf{Rows} & \textbf{Cols} & \textbf{Types} \\
\midrule
Adult Income     & Socioeconomic   & 48,842 & 14  & Mixed \\
Diabetes (PIMA)  & Clinical        & 768    & 9   & Numeric \\
Credit Default   & Financial       & 30,000 & 24  & Mixed \\
\bottomrule
\end{tabular}
\end{table}

Synthetic datasets are generated using CTGAN, TVAE, and GReaT via the
SDV framework~\cite{patki2016synthetic}, with default hyperparameters.

\subsection{LLM Discriminator}

We evaluate two commercial LLMs: GPT-4o (OpenAI) and Claude~3.5~Sonnet
(Anthropic).
Each model receives a fixed system prompt instructing it to behave as a
data scientist evaluating data authenticity, and is required to respond
in a structured JSON schema containing:
(i)~a binary \texttt{verdict} (\texttt{REAL}/\texttt{SYNTHETIC}),
(ii)~a \texttt{confidence} score (0--100),
(iii)~a \texttt{reasoning} string, and
(iv)~lists of \texttt{red\_flags} and \texttt{supporting\_evidence}.
Temperature is set to zero for reproducibility.

\subsection{Human Study}

Human participants are recruited from the research group and related
data science programmes.
The labeling interface (Section~\ref{sec:impl}) presents one table per
trial, collects a binary verdict, a 0--100 confidence slider, and an
optional free-text reasoning field.
Trials are balanced across labels and conditions and randomised per
participant.

\subsection{Evaluation Metrics}

We report per-cell accuracy, AUROC (using confidence scores), and a
one-sided binomial test against the chance baseline ($p = 0.5$) to
assess statistical significance.
Confidence calibration is assessed by binning confidence scores into
quintiles and comparing mean confidence to mean accuracy within each
bin.

% ── 5. Implementation ─────────────────────────────────────────────────────────
\section{Implementation}
\label{sec:impl}

The full experimental framework is implemented in Python~3.10.
LLM API calls are handled via the \texttt{openai} and
\texttt{anthropic} Python SDKs.
Metadata extraction uses \texttt{pandas}, \texttt{numpy}, and
\texttt{scipy.stats}.
The human labeling interface is built with \texttt{Streamlit} and is
deployable on SURF Research Cloud or the VU DAS cluster via SSH port
forwarding.
Results are logged in newline-delimited JSON (\texttt{.jsonl}) for
reproducibility.
Analysis is conducted in a Jupyter notebook using \texttt{scipy},
\texttt{matplotlib}, and \texttt{seaborn}.
The environment is specified in a \texttt{conda} environment file
targeting Python~3.10 for compatibility with SURF/DAS compute nodes.

% ── 6. Expected Results ───────────────────────────────────────────────────────
\section{Expected Results and Hypotheses}
\label{sec:hypotheses}

Based on prior work on LLM evaluation capabilities and human perception
of synthetic data in other modalities, we formulate the following
directional hypotheses:

\begin{description}[leftmargin=0pt,itemsep=4pt]
  \item[$H_1$] Both LLMs and humans will achieve above-chance accuracy
    in condition C1, but accuracy will be modest (55--70\%), reflecting
    the inherent difficulty of visual table inspection.
  \item[$H_2$] Metadata (C2) will significantly increase LLM accuracy,
    as LLMs are well-suited to integrating statistical descriptions.
    The improvement for humans may be smaller or inconsistent,
    due to the cognitive cost of processing distributional information.
  \item[$H_3$] LLMs will be better calibrated than humans:
    their stated confidence will correlate more strongly with actual
    accuracy, as observed in prior LLM-as-judge studies~\cite{zheng2023judging}.
  \item[$H_4$] LLMs will predominantly cite statistical cues
    (skewness anomalies, implausible correlations),
    while humans will cite semantic cues
    (implausible value combinations, contextual implausibility).
  \item[$H_5$] LLM discrimination accuracy will positively
    correlate with standard fidelity metrics: tables that are
    statistically closer to real data (lower TVD, KS) will be
    harder for LLMs to detect as synthetic.
\end{description}

% ── 7. Discussion ─────────────────────────────────────────────────────────────
\section{Limitations and Future Work}
\label{sec:discussion}

Several limitations bound the present study.
First, the set of synthesis methods is restricted to widely used but
ultimately parametric or autoregressive models; diffusion-based
tabular generators~\cite{jolicoeur2023generating} may produce data
with distinct artifact signatures.
Second, the human participant pool is drawn from a data-science-literate
population; results may not generalise to domain experts (e.g.,
clinicians judging medical records) or lay users.
Third, LLM prompts expose only the first $N$ rows; alternative
presentation strategies (random sampling, stratified sampling) may
affect discrimination performance.

Future work will extend the framework to (i)~additional synthesis
methods including REaLTabFormer and diffusion-based generators,
(ii)~domain-expert human cohorts, (iii)~multi-table relational datasets,
and (iv)~an adversarial setting in which synthesis methods are
optimised to fool the LLM discriminator.

% ── 8. Conclusion ─────────────────────────────────────────────────────────────
\section{Conclusion}
\label{sec:conclusion}

We have introduced a \emph{LLM-as-Discriminator} framework for
evaluating the perceptual realism of synthetic tabular data, and
presented a controlled experimental design comparing LLM and human
evaluators across two information conditions.
The proposed framework fills a gap in the tabular data evaluation
landscape: while aggregate fidelity metrics are well-established,
the question of \emph{perceived} realism---whether an informed observer
can detect synthesis---has not been empirically studied.
Our open-source implementation enables replication on standard compute
infrastructure and straightforward extension to new datasets and
synthesis methods.

% ── References ────────────────────────────────────────────────────────────────
% NOTE: Use splncs04 on Overleaf or a full TeX Live install.
% splncs04.bst ships with the official Springer LNCS package from
% https://www.springernature.com/gp/authors/campaigns/latex-author-support
% For local compilation without the full package, plain is used here.
\bibliographystyle{plain}
\bibliography{references}

\end{document}
